\chapter{Clitoral Pain (Clitorodynia)}

\section*{Definition}
Pain, discomfort, hypersensitivity, or altered sensation localised to the clitoris, which may be burning, stabbing, or aching, and can occur spontaneously or with stimulation.

\section*{What Happens in Your Body}
Clitoral tissue is densely innervated with sensory nerve endings and contains abundant oestrogen and androgen receptors. Declining oestrogen and testosterone during menopause cause clitoral tissue atrophy, reduced blood flow, and altered nerve sensitivity. The clitoral nerve is a terminal branch of the pudendal nerve (S2-S4); entrapment or irritation may cause referred or localised pain. Vulvodynia and clitorodynia share overlapping neuropathic and musculoskeletal mechanisms.

\section*{Causes}
\begin{itemize}
  \item Genitourinary syndrome of menopause (clitoral atrophy)
  \item Vulvodynia (generalised or localised vulvar pain)
  \item Pudendal nerve entrapment or neuralgia
  \item Trauma (childbirth, surgery, sexual)
  \item Lichen sclerosus or lichen planus
  \item Recurrent candidiasis or fungal infection
  \item Contact dermatitis (soaps, lubricants, laundry products)
  \item Peyronie disease (affecting clitoral suspensory ligament, rare)
  \item Endometriosis of the clitoris (rare)
  \item Psychosexual factors (anxiety, history of abuse)
\end{itemize}

\section*{When to See a Doctor}
See your doctor if:
\begin{itemize}
  \item Symptoms are severe or getting worse over time
  \item Clitoral pain with visible lesions, discharge, fever, or associated with bladder/bowel symptoms, or pain that interferes with daily function
  \item You have other concerning symptoms such as fever, weight loss, or bleeding
\end{itemize}

\section*{Related Symptoms}
\begin{itemize}
  \item Vulvovaginal irritation
  \item Vaginal dryness
  \item Painful sex
  \item Pelvic pain
  \item Labia changes
\end{itemize}

\textbf{Important:} If this symptom persists, your doctor will check for serious underlying causes.

\section*{Self-Care / Home Management}
\begin{itemize}
  \item Rest and avoid activities that make it worse.
  \item Maintain adequate hydration and a balanced diet.
  \item Over-the-counter remedies may provide symptomatic relief (consult a pharmacist or doctor).
  \item Monitor symptoms and seek professional advice if they persist or worsen.
  \item Avoid self-medicating with prescription medications without medical guidance.
\end{itemize}

\section*{How Your Doctor Will Evaluate You}
\begin{itemize}
  \item A thorough history and physical examination are the first steps in evaluation.
  \item Laboratory tests (blood work, urinalysis) may be ordered based on clinical suspicion.
  \item Imaging studies (X-ray, ultrasound, CT, MRI) may be indicated when structural pathology is suspected.
  \item Specialist referral is arranged when the diagnosis remains uncertain or requires specific expertise.
\end{itemize}

\section*{Treatment Options}
\begin{itemize}
  \item Treatment targets the underlying cause identified through diagnostic evaluation.
  \item Supportive care and symptom management are provided while awaiting definitive therapy.
  \item Medications, lifestyle modifications, and physical therapies may be prescribed as appropriate.
  \item Follow-up ensures treatment effectiveness and allows timely adjustments to the care plan.
\end{itemize}

\clearpage